\numberlesssection{ЗАКЛЮЧЕНИЕ}
В ходе выполнения курсовой работы поставленная цель достигнута --- 
была успешно спроектирована и разработана база данных для анализа 
отказоустойчивости систем на основе микросервисной архитектуры, а также 
реализовано прикладное веб-приложение, предоставляющее пользовательский 
интерфейс для работы с ней.

Все поставленные задачи выполнены:
\begin{itemize}
    \item проведен детальный анализ предметной области обеспечения надежности
    распределенных систем, а также существующих систем мониторинга, подтвердивший
    актуальность разработки инструмента для статического анализа архитектурных уязвимостей;
    \item сформулированы требования к системе и обоснован выбор
    гибридного подхода к хранению информации: использование реляционной СУБД
    PostgreSQL для управления доступом и метаданных и графовой СУБД Neo4j для
    оптимизированной работы с топологией;
    \item разработана ролевая модель управления правами доступа пользователей на
    уровне приложения, обеспечивающая безопасную командную работу и изоляцию проектов;
    \item спроектированы структуры реляционной базы данных и модели графа свойств, а
    корректность реляционной схемы подтверждена успешным прохождением проверки на
    соответствие требованиям третьей нормальной формы;
    \item выбраны средства разработки: язык программирования C\#,
    платформу .NET 10, фреймворк Entity Framework Core и расширение процедур APOC
    для Neo4j;
    \item описана методика и проведено автоматизированное тестирование программного
    обеспечения путем написания модульных тестов с использованием библиотеки Moq и
    интеграционных тестов баз данных в xUnit;
    \item проведено исследование временной эффективности разработанных
    алгоритмов, которое подтвердило высокую производительность и масштабируемость
    графовой базы данных на топологиях, содержащих до десяти тысяч узлов и десятки
    циклических зависимостей.
\end{itemize}
